\bcpchapter{The Order for the Visitation of the Sick}

\begin{rubric}
When any person is sick, notice shall be given thereof to the Minister of the Parish; who, coming into the sick person's house, shall say,
\end{rubric}

\lett{P}{eace} be to this house, and to all that dwell in it.

\begin{rubric}
When he cometh into the sick person's presence he shall say, kneeling down,
\end{rubric}

\lett{R}{emember} not, Lord, our iniquities, nor the iniquities of our forefathers: Spare us, good Lord, spare thy people, whom thou hast redeemed with thy most precious blood, and be not angry with us for ever.

\inrub{Answer.} Spare us, good Lord.

\begin{rubric}
Then the Minister shall say,
\end{rubric}

\lup{Let us pray.}

\V{Lord, have mercy upon us.}
\R{Christ, have mercy upon us.}
\V{Lord, have mercy upon us.}

\lett{O}{ur} Father, which art in heaven, Hallowed be thy Name; Thy kingdom come; Thy will be done; In earth as it is in heaven. Give us this day our daily bread. And forgive us our trespasses, As we forgive them that trespass against us. And lead us not into temptation; But deliver us from evil.

\inlineresp{Amen.}

\V{O Lord, save thy servant;}
\R{Which putteth \changeable{his} trust in thee}
\V{Send \changeable{him} help from thy holy place;}
\R{And evermore mightily defend \changeable{him}.}
\V{Let the enemy have no advantage of \changeable{him};}
\R{Nor the wicked approach to hurt \changeable{him}.}
\V{Be unto \changeable{him}, O Lord, a strong tower.}
\R{From the face of \changeable{his} enemy.}
\V{O Lord, hear our prayers.}
\R{And let our cry come unto thee.}

\newpage

\begin{rubric}
Minister.
\end{rubric}

\lett{O}{ Lord}, look down from heaven, behold, visit, and relieve this thy servant.
Look upon \changeable{him} with the eyes of thy mercy, give \changeable{him} comfort and sure confidence in thee, defend him from the danger of the enemy, and keep \changeable{him} in perpetual peace and safety; through Jesus Christ our Lord.
\inlineresp{Amen.}

\lett{H}{ear} us, Almighty and most merciful God and Saviour; extend thy accustomed goodness to this thy servant who is grieved with sickness.
Sanctify, we beseech thee, this thy fatherly correction to \changeable{him}; that the sense of \changeable{his} weakness may add strength to \changeable{his} faith, and seriousness to \changeable{his} repentance:
That, if it shall be thy good pleasure to restore \changeable{him} to \changeable{his} former health, \changeable{he} may lead the residue of \changeable{his} life in thy fear, and to thy glory:
or else, give \changeable{him} grace so to take thy visitation, that, after this painful life ended, \changeable{he} may dwell with thee in life everlasting; through Jesus Christ our Lord.
\inlineresp{Amen.}

\begin{rubric}
Then shall the Minister exhort the sick person after this form, or other like.
\end{rubric}

\lett{D}{early} beloved, know this, that Almighty God is the Lord of life and death, and of all things to them pertaining, as youth, strength, health, age, weakness, and sickness.
Wherefore, whatsoever your sickness is, know you certainly, that it is God's visitation. And for what cause soever this sickness is sent unto you; whether it be to try your patience for the example of others, and that your faith may be found in the day of the Lord laudable, glorious, and honourable, to the increase of glory and endless felicity;
or else it be sent unto you to correct and amend in you whatsoever doth offend the eyes of your heavenly Father; know you certainly, that if you truly repent you of your sins, and bear your sickness patiently, trusting in God's mercy, for his dear Son Jesus Christ's sake, and render unto him humble thanks for his fatherly visitation, submitting yourself wholly unto his will, it shall turn to your profit, and help you forward in the right way that leadeth unto everlasting life.

\newpage

\begin{rubric}
If the person visited be very sick, then the Curate may end his exhortation in this place, or else proceed.
\end{rubric}

\lett{T}{ake} therefore in good part the chastisement of the Lord:
For (as Saint Paul saith in the twelfth Chapter to the Hebrews) whom the Lord loveth he chasteneth, and scourgeth every son whom he receiveth.
If ye endure chastening, God dealeth with you as with sons; for what son is he whom the father chasteneth not?
But if ye be without chastisement, whereof all are partakers, then are ye bastards, and not sons.
Furthermore, we have had fathers of our flesh, which corrected us, and we gave them reverence: shall we not much rather be in subjection unto the Father of spirits, and live?
For they verily for a few days chastened us after their own pleasure; but he for our profit, that we might be partakers of his holiness.
These words, good brother, are written in holy Scripture for our comfort and instruction; that we should patiently, and with thanksgiving, bear our heavenly Father's correction, whensoever by any manner of adversity it shall please his gracious goodness to visit us.
And there should be no greater comfort to Christian persons, than to be made like unto Christ, by suffering patiently adversities, troubles, and sicknesses.
For he himself went not up to joy, but first he suffered pain; he entered not into his glory before he was crucified.
So truly our way to eternal joy is to suffer here with Christ; and our door to enter into eternal life is gladly to die with Christ; that we may rise again from death, and dwell with him in everlasting life.
Now therefore, taking your sickness, which is thus profitable for you, patiently, I exhort you, in the Name of God, to remember the profession which you made unto God in your Baptism.
And forasmuch as after this life there is an account to be given unto the righteous judge, by whom all must be judged, without respect of persons, I require you to examine yourself and your estate, both toward God and man; so that, accusing and condemning yourself for your own faults, you may find mercy at our heavenly Father's hand for Christ's sake, and not be accused and condemned in that fearful judgement.
Therefore I shall rehearse to you the Articles of our Faith, that you may know whether you do believe as a Christian man should, or no.

\newpage

\begin{rubric}
Here the Minister shall rehearse the Articles of the Faith, saying thus,
\end{rubric}

\lett{D}{ost} thou believe in God the Father Almighty, Maker of heaven and earth?

And in Jesus Christ his only-begotten Son our Lord? And that he was conceived by the Holy Ghost, born of the Virgin Mary; that he suffered under Pontius Pilate, was crucified, dead, and buried; that he went down into hell, and also did rise again the third day; that he ascended into heaven, and sitteth at the right hand of God the Father Almighty; and from thence shall come again at the end of the world, to judge the quick and the dead?

And dost thou believe in the Holy Ghost; the holy Catholick Church; the Communion of Saints; the Remission of sins; the Resurrection of the flesh; and everlasting life after death?

\begin{rubric}
The sick person shall answer,
\end{rubric}

All this I stedfastly believe.

\begin{rubric}
Then shall the Minister examine whether he repent him truly of his sins, and be in charity with all the world; exhorting him to forgive, from the bottom of his heart, all persons that have offended him; and if he hath offended any other, to ask them forgiveness; and where he hath done injury or wrong to any man, that he make amends to the uttermost of his power.
And if he have not before disposed of his goods, let him then be admonished to make his Will, and to declare his debts, what he oweth, and what is owing unto him; for the better discharging of his conscience, and the quietness of his Executors.
But men should often be put in remembrance to take order for the settling of their temporal estates, whilst they are in health.
\end{rubric}

\begin{rubric}
These words before rehearsed may be said before the Minister begin his Prayer, as he shall see cause.
\end{rubric}

\begin{rubric}
The Minister should not omit earnestly to move such sick persons as are of ability to be liberal to the poor.
\end{rubric}

\newpage

\begin{rubric}
Here shall the sick person be moved to make a special confession of his sins, if he feel his conscience troubled with any weighty matter.
After which confession, the Priest shall absolve him (if he humbly and heartily desire it) after this sort.
\end{rubric}

\lett{O}{ur} Lord Jesus Christ, who hath left power to his Church to absolve all sinners who truly repent and believe in him, of his great mercy forgive thee thine offences: And by his authority committed to me, I absolve thee from all thy sins, In the Name of the Father, and of the Son, and of the Holy Ghost.
\inlineresp{Amen.}

\begin{rubric}
And then the Priest shall say the Collect following.
\end{rubric}

\lup{Let us pray.}

\lett{O}{ Most} merciful God, who, according to the multitude of thy mercies, dost so put away the sins of those who truly repent, that thou rememberest them no more: Open thine eye of mercy upon this thy servant, who most earnestly desireth pardon and forgiveness.
Renew in \changeable{him}, most loving Father, whatsoever hath been decayed by the fraud and malice of the devil, or by \changeable{his} own carnal will and frailness; preserve and continue this sick member in the unity of the Church; consider \changeable{his} contrition, accept his tears, asswage his pain, as shall seem to thee most expedient for \changeable{him}.
And forasmuch as \changeable{he} putteth \changeable{his} full trust only in thy mercy, impute not unto him \changeable{his} former sins, but strengthen \changeable{him} with thy blessed Spirit; and, when thou art pleased to take \changeable{him} hence, take \changeable{him} unto thy favour, through the merits of thy most dearly beloved Son Jesus Christ our Lord.
\inlineresp{Amen.}

\begin{rubric}
Then shall the Minister say this Psalm,
\end{rubric}

\subsection{Psalm 71\\
(In te, Domine, speravi)}
In thee, O Lord, have I put my trust, let me never be put to confusion : but rid me and deliver me in thy righteousness, incline thine ear unto me, and save me.\\
Be thou my strong hold, whereunto I may alway resort : thou hast promised to help me, for thou art my house of defence and my castle.\\
Deliver me, O my God, out of the hand of the ungodly : out of the hand of the unrighteous and cruel man.\\
For thou, O Lord God, art the thing that I long for : thou art my hope, even from my youth.\\
\newpage
Through thee have I been holden up ever since I was born : thou art he that took me out of my mother's womb; my praise shall be always of thee.\\
I am become as it were a monster unto many : but my sure trust is in thee.\\
O let my mouth be filled with thy praise : that I may sing of thy glory and honour all the day long.\\
Cast me not away in the time of age : forsake me not when my strength faileth me.\\
For mine enemies speak against me, and they that lay wait for my soul take their counsel together, saying : God hath forsaken him; persecute him, and take him, for there is none to deliver him.\\
Go not far from me, O God : my God, haste thee to help me.\\
Let them be confounded and perish that are against my soul : let them be covered with shame and dishonour that seek to do me evil.\\
As for me, I will patiently abide alway : and will praise thee more and more.\\
My mouth shall daily speak of thy righteousness and salvation : for I know no end thereof.\\
I will go forth in the strength of the Lord God : and will make mention of thy righteousness only.\\
Thou, O God, hast taught me from my youth up until now : therefore will I tell of thy wondrous works.\\
Forsake me not, O God, in mine old age, when I am gray-headed : until I have shewed thy strength unto this generation, and thy power to all them that are yet for to come.\\
Thy righteousness, O God, is very high : and great things are they that thou hast done; O God, who is like unto thee?\\
O what great troubles and adversities hast thou shewed me, and yet didst thou turn and refresh me : yea, and broughtest me from the deep of the earth again.\\
Thou hast brought me to great honour : and comforted me on every side.\\
Therefore will I praise thee and thy faithfulness, O God, playing upon an instrument of musick : unto thee will I sing upon the harp, O thou Holy One of Israel.\\
My lips will be fain when I sing unto thee : and so will my soul whom thou hast delivered.\\
My tongue also shall talk of thy righteousness all the day long : for they are confounded and brought unto shame that seek to do me evil.\\
\newpage
Glory be to the Father and to the Son : and to the Holy Ghost;\\
As it was in the beginning, is now, and ever shall be : world without end. Amen.

\begin{rubric}
Adding this.
\end{rubric}

\lett{O}{ Saviour} of the world, who by thy Cross and precious Blood hast redeemed us, Save us, and help us, we humbly beseech thee, O Lord.

\begin{rubric}
Then shall the Minister say,
\end{rubric}

\lett{T}{he} Almighty Lord, who is a most strong tower to all them that put their trust in him, to whom all things in heaven, in earth, and under the earth, do bow and obey, be now and evermore thy defence; and make thee know and feel, that there is none other Name under heaven given to man, in whom, and through whom, thou mayest receive health and salvation, but only the Name of our Lord Jesus Christ.
\inlineresp{Amen.}

\begin{rubric}
And after that shall say,
\end{rubric}

\lett{U}{nto} God's gracious mercy and protection we commit thee. The Lord bless thee, and keep thee. The Lord make his face to shine upon thee, and be gracious unto thee. The Lord lift up his countenance upon thee, and give thee peace, both now and evermore.
\inlineresp{Amen.}

\begin{rubric}
A Prayer for a sick Child.
\end{rubric}

\lett{O}{ Almighty} God, and merciful Father, to whom alone belong the issues of life and death:
Look down from heaven, we humbly beseech thee, with the eyes of mercy upon this child now lying upon the bed of sickness:
Visit \changeable{him}, O Lord, with thy salvation; deliver \changeable{him} in thy good appointed time from \changeable{his} bodily pain, and save \changeable{his} soul for thy mercies' sake:
That, if it shall be thy pleasure to prolong \changeable{his} days here on earth, \changeable{he} may live to thee, and be an instrument of thy glory, by serving thee faithfully, and doing good in \changeable{his} generation;
or else receive \changeable{him} into those heavenly habitations, where the souls of them that sleep in the Lord Jesus enjoy perpetual rest and felicity.
Grant this, O Lord, for thy mercies' sake, in the same thy Son our Lord Jesus Christ, who liveth and reigneth with thee and the Holy Ghost, ever one God, world without end.
\inlineresp{Amen.}

\newpage

\begin{rubric}
A Prayer for a sick person, when there appeareth small hope of recovery.
\end{rubric}

\lett{O}{ Father} of mercies, and God of all comfort, our only help in time of need: We fly unto thee for succour in behalf of this thy servant, here lying under thy hand in great weakness of body.
Look graciously upon \changeable{him}, O Lord; and the more the outward man decayeth, strengthen \changeable{him}, we beseech thee, so much the more continually with thy grace and Holy Spirit in the inner man.
Give \changeable{him} unfeigned repentance for all the errors of \changeable{his} life past, and stedfast faith in thy Son Jesus; that \changeable{his} sins may be done away by thy mercy, and \changeable{his} pardon sealed in heaven, before \changeable{he} go hence, and be no more seen.
We know, O Lord, that there is no word impossible with thee; and that, if thou wilt, thou canst even yet raise \changeable{him} up, and grant \changeable{him} a longer continuance amongst us: Yet, forasmuch as in all appearance the time of \changeable{his} dissolution draweth near, so fit and prepare \changeable{him}, we beseech thee, against the hour of death, that after \changeable{his} departure hence in peace, and in thy favour, \changeable{his} soul may be received into thine everlasting kingdom, through the merits and mediation of Jesus Christ, thine only Son, our Lord and Saviour.
\inlineresp{Amen.}

\begin{rubric}
A commendatory Prayer for a sick person at the point of departure.
\end{rubric}

\lett{O}{ Almighty} God, with whom do live the spirits of just men made perfect, after they are delivered from their earthly prisons: We humbly commend the soul of this thy servant, our dear \changeable{brother}, into thy hands, as into the hands of a faithful Creator, and most merciful Saviour; most humbly beseeching thee, that it may be precious in thy sight.
Wash it, we pray thee, in the blood of that immaculate Lamb, that was slain to take away the sins of the world; that whatsoever defilements it may have contracted in the midst of this miserable and naughty world, through the lusts of the flesh, or the wiles of Satan, being purged and done away, it may be presented pure and without spot before thee.
And teach us who survive, in this and other like daily spectacles of mortality, to see how frail and uncertain our own condition is; and so to number our days, that we may seriously apply our hearts to that holy and heavenly wisdom, whilst we live here, which may in the end bring us to life everlasting, through the merits of Jesus Christ thine only Son our Lord.
\inlineresp{Amen.}

\newpage

\begin{rubric}
A Prayer for persons troubled in mind or in conscience.
\end{rubric}

\lett{O}{ Blessed} Lord, the Father of mercies, and the God of all comforts: We beseech thee, took down in pity and compassion upon this thy afflicted servant.
Thou writest bitter things against \changeable{him}, and makest \changeable{him} to possess his former iniquities; thy wrath lieth hard upon \changeable{him}, and \changeable{his} soul is full of trouble:
But, O merciful God, who hast written thy holy Word for our learning, that we, through patience and comfort of thy holy Scriptures, might have hope;
give \changeable{him} a right understanding of himself, and of thy threats and promises; that \changeable{he} may neither cast away \changeable{his} confidence in thee, nor place it any where but in thee.
Give \changeable{him} strength against all \changeable{his} temptations, and heal all his distempers.
Break not the bruised reed, nor quench the smoking flax.
Shut not up thy tender mercies in displeasure; but make \changeable{him} to hear of joy and gladness, that the bones which thou hast broken may rejoice.
Deliver \changeable{him} from fear of the enemy, and lift up the light of thy countenance upon \changeable{him}, and give \changeable{him} peace, through the merits and mediation of Jesus Christ our Lord.
\inlineresp{Amen.}
